\documentclass[11pt,a4paper]{article}
\usepackage[T1]{fontenc}
\usepackage{lmodern}
\usepackage[margin=1in]{geometry}
\usepackage{amsmath,amssymb,amsthm,mathtools}
\usepackage{microtype}
\usepackage{booktabs,longtable,array}
\usepackage{xcolor}
\usepackage{listings}
\usepackage[colorlinks=true,linkcolor=blue,citecolor=blue,urlcolor=blue]{hyperref}
\hypersetup{pdftitle={A Proof of Colton's Conjecture on Refactorable Numbers},pdfauthor={Keio}}

\newtheorem{theorem}{Theorem}[section]
\newtheorem{proposition}[theorem]{Proposition}
\newtheorem{lemma}[theorem]{Lemma}
\newtheorem{corollary}[theorem]{Corollary}
\newtheorem{conjecture}[theorem]{Conjecture}
\theoremstyle{definition}
\newtheorem{definition}[theorem]{Definition}
\newtheorem{remark}[theorem]{Remark}

\newcommand{\cS}{\mathcal S}
\newcommand{\cU}{\mathcal U}
\newcommand{\cT}{\mathcal T}
\newcommand{\N}{\mathbb N}

\lstset{basicstyle=\ttfamily\scriptsize,breaklines=true,columns=fullflexible,keepspaces=true,
  frame=single,showstringspaces=false,upquote=true,tabsize=4}

\title{\bfseries A Proof of Colton's Conjecture on Refactorable Numbers\\[6pt]
\large Smooth--rough $\tau$-decomposition, certified checkpoint chains, and a template tail}
\author{Keio}
\date{24 October 2026}

\begin{document}
\maketitle

\begin{abstract}
A positive integer $k$ is \emph{refactorable} (a \emph{tau number}) if the number of its divisors $\tau(k)$ divides $k$. Let $T(n)$ be the number of refactorable integers in $[1,n]$. In 1999 Colton conjectured that $T(n)\ge \pi(n)/2$ for every $n$. In 2002 Zelinsky proved the inequality for all sufficiently large $n$, with the explicit bound $7.42\cdot10^{13}$ for possible counterexamples. The full statement has remained open and is recorded as open in the Formal Conjectures repository.

We prove the conjecture: $\pi(n)\le 2T(n)$ for every $n\ge1$, with equality only at $n=7$. The proof does not use Zelinsky's bound. It rests on one observation: for $k<2^{52}$, every prime factor of $\tau(k)$ is at most $47$. Hence refactorability is decided by the $47$-smooth part of $k$ together with the exponent pattern of its $47$-rough part. This \emph{smooth--rough $\tau$-decomposition} gives three things:
\begin{enumerate}
\item[(i)] an exact formula expressing $T(x)$ for $x<2^{52}$ through values of the prime-counting function, which can be evaluated in sublinear time;
\item[(ii)] a certified lower bound $L(x)\le T(x)$ for all $x$;
\item[(iii)] finite families of \emph{templates}, integers whose products with large primes are always refactorable.
\end{enumerate}
Combined with the Rosser--Schoenfeld bounds, (iii) settles every $x\ge10^{14}$ through one explicit finite inequality. A monotone checkpoint principle turns finitely many values of $L$ and $\pi$ into proofs on whole intervals; $27$ certified links close the range $[10^{11},10^{14}]$. The range $n\le10^{11}$ is covered by an exhaustive sieve, and independently by a $58$-link chain starting at $10^{6}$.

As a by-product we obtain three new terms of OEIS A111291:
\[
T(10^{12})=26\,568\,322\,474,\quad T(10^{13})=251\,342\,379\,436,\quad T(10^{14})=2\,388\,671\,712\,980.
\]
Each was confirmed by several different decompositions. All programs are listed in full in the appendix.
\end{abstract}

\tableofcontents

\section{Introduction}

\subsection{Refactorable numbers and Colton's conjecture}

Let $\tau(k)$ denote the number of positive divisors of $k$. Kennedy and Cooper~\cite{KC90} called $k$ a \emph{tau number} if $\tau(k)\mid k$, and showed that these integers have natural density zero. Colton~\cite{Co99} rediscovered them with his automated theory-formation program HR and called them \emph{refactorable}. They form OEIS sequence A033950~\cite{OEIS}:
\[
1,\ 2,\ 8,\ 9,\ 12,\ 18,\ 24,\ 36,\ 40,\ 56,\ 60,\ 72,\ 80,\ 84,\ 88,\ 96,\ \dots
\]
Throughout we write
\[
T(x)=\#\{1\le k\le x:\ \tau(k)\mid k\},\qquad \pi(x)=\#\{p\le x:\ p \text{ prime}\}\qquad(x\ge 0).
\]
Spiro~\cite{Sp85} determined the order of magnitude of $T$: by his asymptotic formula, as recorded in~\cite{FC},
\[
T(x)=\frac{x}{\sqrt{\log x}}\,(\log\log x)^{-1+o(1)}.
\]
In particular $T(x)/\pi(x)\to\infty$. The convergence is very slow, however: $2T(x)/\pi(x)$ is $1.095$ at $x=10^{3}$ and only $1.491$ at $x=10^{14}$ (Table~\ref{tab:values}).

Colton~\cite{Co99} conjectured that refactorable numbers are never less than half as frequent as the primes. We use the pointwise formulation stated by Zelinsky~\cite{Ze02} and recorded in~\cite{FC}.

\begin{conjecture}[Colton]\label{conj:colton}
For every positive integer $n$, $T(n)\ge\frac12\pi(n)$; equivalently, $\pi(n)\le 2T(n)$.
\end{conjecture}

Zelinsky~\cite{Ze02} proved that for every fixed $c>0$ one has $T(n)>c\,\pi(n)$ for all sufficiently large $n$. He also gave the explicit bound $7.42\cdot10^{13}$, beyond which Conjecture~\ref{conj:colton} holds. The finite range below this bound was left open. The Formal Conjectures repository~\cite{FC} (snapshot of 23 October 2026) states the conjecture in Lean as
\begin{center}
\texttt{theorem colton\_conjecture :\ $\forall$ n :\ $\N$, Nat.primeCounting n $\le$ 2 * countRefactorableNat n}
\end{center}
It is tagged \texttt{research open}, with the remark that Zelinsky proved it beyond $7.42\cdot10^{13}$ ``and the remaining range is open''. There \texttt{countRefactorableNat n} counts the $k\in[1,n]$ with $\tau(k)\mid k$, so it is our $T(n)$.

A weaker asymptotic formulation, $T(x)\ge x/(2\log x)$ for all sufficiently large $x$, is recorded in~\cite{FC} as solved, as a consequence of~\cite{Ze02,Sp85}.

Verifying the remaining range by brute force means computing $\tau(k)$ for every $k<7.42\cdot10^{13}$. Our parallel sieve needs about $40$ minutes on four cores just to reach $10^{11}$, so this is a substantial computation (Section~\ref{sec:timing}).

\subsection{Main result}

\begin{theorem}\label{thm:main}
For every integer $n\ge1$,
\[
\pi(n)\ \le\ 2\,T(n),
\]
with equality if and only if $n=7$. In particular Conjecture~\ref{conj:colton} is true. The statement \texttt{colton\_conjecture} of~\cite{FC} also holds; for $n=0$ both sides vanish.
\end{theorem}

\begin{corollary}[Explicit asymptotic form]\label{cor:explicit}
$T(x)>x/(2\log x)$ for every real $x\ge2$, and this inequality fails for every $x\in(1,2)$.
\end{corollary}

\begin{proof}
For $x\ge17$, Theorem~\ref{thm:main} and Lemma~\ref{lem:RS}(a) give
\[
T(x)=T(\lfloor x\rfloor)\ge\tfrac12\pi(\lfloor x\rfloor)=\tfrac12\pi(x)>\frac{x}{2\log x}.
\]
Let $f(x)=x/(2\log x)$. On $[2,17)$ the function $T$ equals $2,3,4,5$ on the intervals $[2,8)$, $[8,9)$, $[9,12)$, $[12,17)$ respectively. Since $f$ decreases on $(1,e]$ and increases on $[e,\infty)$, on each of these intervals $f$ is bounded by its values at the two endpoints. These values are $f(2)<1.45$, $f(8)<1.93$, $f(9)<2.05$, $f(12)<2.42$ and $f(17)<3.01$, each smaller than the corresponding value of $T$. Finally, for $1<x<2$ we have $2\log x<x$, so $f(x)>1=T(x)$.
\end{proof}

At $n=7$ we have $\pi(7)=4$ and $T(7)=2$, the refactorable numbers up to $7$ being $1$ and $2$. For $100\le n\le10^{8}$ the ratio $2T(n)/\pi(n)$ is smallest at $n=439$, where it equals $2\cdot44/85=88/85\approx1.0353$ (Section~\ref{sec:exhaustive}).

\subsection{The method}

The proof does not use Zelinsky's bound. It starts from an elementary observation (Lemma~\ref{lem:threshold}):
\begin{quote}
\emph{If $k<2^{52}$, then every prime factor of $\tau(k)$ is at most $47$.}
\end{quote}
Indeed, every exponent $e$ in the factorization of such a $k$ satisfies $e\le51$, so every factor $e+1$ of $\tau(k)$ is at most $52<53$. Write $k=s\,r$, where $s$ is the $47$-smooth part and $r$ the $47$-rough part of $k$. Then $\tau(k)$ is coprime to $r$, and therefore
\[
\tau(k)\mid k\iff \tau(s)\,\tau(r)\mid s .
\]
So the rough part matters only through its exponent pattern. When $r$ is squarefree, only the number $\omega(r)$ of its prime factors matters. We call this the \emph{smooth--rough $\tau$-decomposition}. It turns a divisibility condition on $k$ into a counting problem for almost-primes, which reduces to prime counting. We use it in three ways.

\begin{enumerate}
\item[(A)] \textbf{Exact formula} (Theorem~\ref{thm:exact}). For $x<2^{52}$, $T(x)$ is an explicit finite combination of counting functions $\Phi^{Q}_j$ of squarefree $47$-rough integers with $j$ prime factors. By Lemma~\ref{lem:phirec} these reduce to values of $\pi$. Combinatorial prime-counting algorithms~\cite{DR96,Go01,primecount} then compute $T(10^{14})$ in under three minutes.
\item[(B)] \textbf{Certified checkpoint chains} (Section~\ref{sec:chains}). The decomposition also gives, for \emph{all} $x$, a lower bound $L(x)\le T(x)$ (Proposition~\ref{prop:L}). Since $\pi$ and $T$ are nondecreasing, a single inequality $\pi(b)<2L(a)$ proves $\pi(n)<2T(n)$ for \emph{every} $n\in[a,b]$ (Lemma~\ref{lem:checkpoint}). Twenty-seven such links cover $[10^{11},10^{14}]$.
\item[(C)] \textbf{Template tail} (Theorem~\ref{thm:tail}). A \emph{template} is an integer $m=s\,p_1\cdots p_{j-1}$ such that $mq$ is refactorable for every prime $q>p_{j-1}$. A finite family of about $2\cdot10^7$ templates, combined with the Rosser--Schoenfeld inequalities~\cite{RS62}, gives
\[
T(x)>0.53934\,\frac{x}{\log x}-2.41\cdot 10^{10}\qquad(x\ge10^{14}).
\]
This implies $\pi(x)<2T(x)$ for every real $x\ge10^{14}$.
\end{enumerate}

The work follows the scheme \emph{new problem $\to$ new method $\to$ application}. The new problem is the following.
\begin{quote}
\textbf{Prime-count reduction problem for $\tau$-divisibility.} Can the counting function $T(x)$, defined by a divisor-function condition, be computed exactly, and bounded below with certificates, using only values of the prime-counting function? On which range of $x$?
\end{quote}
Theorem~\ref{thm:exact} solves the exact version for all $x<2^{52}$; the threshold $2^{52}$ is sharp for the underlying smoothness property (Lemma~\ref{lem:threshold}). Proposition~\ref{prop:L} and Lemma~\ref{lem:template} solve the lower-bound version for all $x$. Theorem~\ref{thm:main} is the application to a known open problem. We are not aware of an earlier use of this reduction to compute or bound $T(x)$.

\subsection{New numerical data}

Theorem~\ref{thm:exact} yields three new terms of OEIS A111291, the number of refactorable numbers $\le10^n$:
\[
a(12)=26\,568\,322\,474,\qquad a(13)=251\,342\,379\,436,\qquad a(14)=2\,388\,671\,712\,980 .
\]
Each value was recomputed with every applicable smoothness bound $B\in\{37,43,47,53\}$, and all runs agree (Section~\ref{sec:exactdata}). The same program reproduces all previously known terms $a(0),\dots,a(11)$.

\subsection{Computer assistance and verifiability}

The proof is computer-assisted. It uses three computations:
\begin{enumerate}
\item[(a)] an exhaustive sieve up to $10^{11}$;
\item[(b)] $85$ evaluations each of $L$ and $\pi$ along two chains;
\item[(c)] one finite sum over about $2\cdot10^7$ templates.
\end{enumerate}
Every number used by a chain is printed in Appendix~\ref{app:chains}, and the complete source code is in Appendix~\ref{app:code}. Section~\ref{sec:crosscheck} lists the independent cross-checks, among them:
\begin{itemize}
\item a second prime-counting algorithm for all values of $\pi$ used by the chains;
\item a brute-force recount of $L$;
\item two independent implementations of the tail certificate;
\item agreement with the twelve previously known terms $a(0),\dots,a(11)$ of A111291.
\end{itemize}

\section{Notation and classical inputs}

Throughout, $p,q$ denote primes, $\log$ is the natural logarithm, $\omega(k)$ is the number of distinct prime factors of $k$, and $v_p(k)$ is the exponent of $p$ in $k$. A positive integer is \emph{$B$-smooth} if all its prime factors are $\le B$, and \emph{$B$-rough} if all of them are $>B$; the integer $1$ is both. For $k\ge1$ put
\[
s(k)=\prod_{p\le 47}p^{v_p(k)},\qquad r(k)=\frac{k}{s(k)},
\]
the $47$-smooth part and the $47$-rough part of $k$. Note that $\pi(47)=15$ and that $53$ is the least prime exceeding $47$. The functions $\pi$ and $T$ are nondecreasing on $[0,\infty)$.

The only analytic inputs are two classical explicit bounds.

\begin{lemma}[Rosser--Schoenfeld~\cite{RS62}]\label{lem:RS}
\begin{enumerate}
\item[(a)] $\pi(x)>x/\log x$ for all $x\ge17$ \textup{\cite[Corollary~1, (3.5)]{RS62}}.
\item[(b)] $\displaystyle\pi(x)<\frac{x}{\log x}\Bigl(1+\frac{3}{2\log x}\Bigr)$ for all $x>1$ \textup{\cite[Theorem~1, (3.2)]{RS62}}.
\end{enumerate}
\end{lemma}

As a sanity check, both inequalities were also confirmed numerically at all $87$ distinct arguments of $\pi$ that occur in the chains of Section~\ref{sec:chains}.

\section{The smooth--rough \texorpdfstring{$\tau$}{tau}-decomposition}\label{sec:decomp}

\begin{lemma}[$\tau$-smoothness threshold]\label{lem:threshold}
Let $P$ be a prime and $k$ a positive integer with $k<2^{P-1}$. Then every prime factor of $\tau(k)$ is smaller than $P$. The bound is sharp, since $\tau(2^{P-1})=P$. In particular, if $k<2^{52}$ then $\tau(k)$ is $47$-smooth.
\end{lemma}

\begin{proof}
We have $\tau(k)=\prod_{p\mid k}(v_p(k)+1)$. For each prime $p\mid k$, $2^{v_p(k)}\le p^{v_p(k)}\le k<2^{P-1}$, hence $v_p(k)+1\le P-1$. Every prime factor of an integer in $[1,P-1]$ is smaller than $P$. For $P=53$, the primes smaller than $P$ are exactly the primes $\le47$.
\end{proof}

\begin{corollary}\label{cor:criterion}
Let $1\le k<2^{52}$, $s=s(k)$ and $r=r(k)$. Then $\tau(k)\mid k$ if and only if $\tau(s)\,\tau(r)\mid s$.
\end{corollary}

\begin{proof}
Since $\gcd(s,r)=1$, we have $\tau(k)=\tau(s)\tau(r)$. By Lemma~\ref{lem:threshold}, $\tau(k)$ is $47$-smooth and hence coprime to the $47$-rough integer $r$. Therefore $\tau(k)\mid sr$ if and only if $\tau(k)\mid s$.
\end{proof}

We split the rough part further: $r=u\,v$ with
\[
u=\prod_{q\mid r,\ v_q(r)\ge2}q^{v_q(r)},\qquad v=\prod_{q\mid r,\ v_q(r)=1}q .
\]
Thus $u$ is powerful, $v$ is squarefree, $\gcd(u,v)=1$, and $\tau(r)=\tau(u)\,2^{\omega(v)}$. We use the following notation.
\begin{itemize}
\item $\cS=\{s\ge1:\ s \text{ is $47$-smooth and } \tau(s)\mid s\}$. For $s\in\cS$ put $d(s)=s/\tau(s)$ and $J(s)=v_2(d(s))$.
\item $\cU$ is the set of $47$-rough powerful integers, including $1$. For $u\in\cU$, $Q(u)$ is the set of prime factors of $u$.
\item For a finite set $Q$ of primes $>47$, an integer $j\ge0$ and a real $y\ge0$,
\[
\Phi^{Q}_j(y)=\#\Bigl\{v\le y:\ v \text{ squarefree and $47$-rough},\ \omega(v)=j,\ \gcd\Bigl(v,\prod_{q\in Q}q\Bigr)=1\Bigr\},
\]
and $\Phi_j=\Phi_j^{\emptyset}$.
\end{itemize}

\begin{theorem}[Exact formula]\label{thm:exact}
For every real $x$ with $1\le x<2^{52}$,
\begin{equation}\label{eq:exact}
T(x)=\sum_{s\in\cS}\ \sum_{\substack{u\in\cU\\ \tau(u)\mid d(s)}}\ \sum_{j=0}^{v_2(d(s)/\tau(u))}\Phi_j^{Q(u)}\Bigl(\frac{x}{su}\Bigr),
\end{equation}
where only the finitely many pairs $(s,u)$ with $su\le x$ contribute.
\end{theorem}

\begin{proof}
Every integer $k\in[1,x]$ factors uniquely as $k=s\,u\,v$ with $s=s(k)$ $47$-smooth, $u\in\cU$, and $v$ squarefree, $47$-rough and coprime to $u$. Conversely, each such triple with $suv\le x$ comes from exactly one $k$. By Corollary~\ref{cor:criterion}, $k$ is refactorable if and only if $\tau(s)\tau(u)2^{\omega(v)}\mid s$. This holds if and only if
\[
\tau(s)\mid s,\qquad \tau(u)\mid d(s),\qquad 2^{\omega(v)}\mid d(s)/\tau(u),
\]
that is, if and only if $s\in\cS$, $\tau(u)\mid d(s)$ and $\omega(v)\le v_2(d(s)/\tau(u))$. For fixed $(s,u)$, the number of admissible $v$ with $\omega(v)=j$ is $\Phi^{Q(u)}_j(x/(su))$.
\end{proof}

\begin{lemma}[Reduction to prime counts]\label{lem:phirec}
Let $y\ge0$ and let $Q$ be a finite set of primes $>47$. For a prime $p\ge 47$, let $\Phi^{(>p)}_j$ be defined like $\Phi_j$ but counting only integers all of whose prime factors exceed $p$; thus $\Phi_j=\Phi_j^{(>47)}$.
\begin{enumerate}
\item[(i)] $\Phi^{Q}_0(y)=1$ if $y\ge1$ and $0$ otherwise. Moreover $\Phi^{(>p)}_1(y)=\max\{0,\pi(y)-\pi(p)\}$; in particular $\Phi_1(y)=\max\{0,\pi(y)-15\}$.
\item[(ii)] For $j\ge2$,
\[
\Phi^{(>p)}_j(y)=\sum_{p'>p,\ p'^{\,j}\le y}\Phi^{(>p')}_{j-1}(y/p'),
\]
where $p'$ runs over primes.
\item[(iii)] For a prime $q\notin Q$ with $q>47$, and for $j\ge1$,
\[
\Phi^{Q\cup\{q\}}_j(y)=\Phi^{Q}_j(y)-\Phi^{Q\cup\{q\}}_{j-1}(y/q).
\]
\end{enumerate}
\end{lemma}

\begin{proof}
Part (i) is immediate. For (ii), classify $v$ by its least prime factor $p'$. Then $v=p'v'$, where $v'\le y/p'$ has $j-1$ prime factors, all larger than $p'$. Since $v'>p'^{\,j-1}$, such $v$ exist only if $p'^{\,j}<y$. For (iii), the integers counted by $\Phi_j^{Q}(y)$ but divisible by $q$ are exactly $v=qv'$, where $v'\le y/q$ is coprime to $Q\cup\{q\}$ and $\omega(v')=j-1$.
\end{proof}

Hence~\eqref{eq:exact} expresses $T(x)$ entirely through finitely many values of $\pi$: values at points $x/n$ with $n$ a positive integer, and values $\pi(p)$ at primes $p\le\sqrt x$. The largest argument needed is $x/8$, because $s=8$ is the smallest element of $\cS$ with $J(s)\ge1$.

\begin{remark}[Cost]\label{rem:cost}
The work is dominated by the evaluations of $\Phi_2(y)$ for $y=x/(su)$, each of which needs $O(\sqrt y)$ look-ups of $\pi$. The series $\sum_{s\in\cS}s^{-1/2}$ converges quickly because $\cS$ consists of $47$-smooth numbers. So, heuristically, the number of look-ups is $x^{1/2+o(1)}$. To this one adds the cost of the prime-counting oracle: a sieve table up to some bound $Z_0$, and the combinatorial algorithms of Deléglise--Rivat and Gourdon~\cite{DR96,Go01} for the few arguments above $Z_0$. The proof needs no complexity statement; measured running times are in Table~\ref{tab:timing}.
\end{remark}

\begin{proposition}[Universal lower bound]\label{prop:L}
For every real $x\ge1$, with no upper restriction on $x$,
\[
T(x)\ \ge\ L(x):=\sum_{s\in\cS}\ \sum_{j=0}^{J(s)}\Phi_j\Bigl(\frac{x}{s}\Bigr).
\]
For $x<2^{52}$, $L(x)$ is exactly the number of refactorable $k\le x$ whose $47$-rough part is squarefree; these are the terms with $u=1$ in~\eqref{eq:exact}.
\end{proposition}

\begin{proof}
Let $s\in\cS$, $0\le j\le J(s)$, and let $v\le x/s$ be squarefree and $47$-rough with $\omega(v)=j$. Put $k=sv$. Then $\tau(k)=\tau(s)\,2^{j}$. Since $2^j\mid d(s)$, the number $\tau(s)2^j$ divides $\tau(s)d(s)=s$, which divides $k$. Moreover $(s,v)=(s(k),r(k))$, so distinct pairs give distinct $k$. This argument does not use Lemma~\ref{lem:threshold}. The last statement is Theorem~\ref{thm:exact} restricted to $u=1$.
\end{proof}

The bound is very tight. At $x=10^{12}$ we have $L(x)=26\,549\,555\,866$ and $T(x)=26\,568\,322\,474$, so $L/T=0.99929$.

\begin{lemma}[Templates]\label{lem:template}
A \emph{template} consists of:
\begin{itemize}
\item an integer $j\ge1$;
\item an element $s\in\cS$ with $J(s)\ge j$, i.e.\ with $2^j\tau(s)\mid s$;
\item primes $53\le p_1<\dots<p_{j-1}$.
\end{itemize}
Its \emph{value} is $m=s\,p_1\cdots p_{j-1}$. Put $p(m)=p_{j-1}$, and $p(m)=47$ if $j=1$. The value determines the template, since $s=s(m)$ and the $p_i$ are the prime factors of $r(m)$; we therefore identify a template with its value $m$.

For every template $m$ and every prime $q>p(m)$, the integer $mq$ is refactorable, and distinct pairs $(m,q)$ give distinct integers. Consequently, for every finite set $\cT$ of templates and every real $x>0$,
\begin{equation}\label{eq:template}
T(x)\ \ge\ \sum_{m\in\cT}\Bigl(\pi\Bigl(\frac{x}{m}\Bigr)-\pi\bigl(p(m)\bigr)\Bigr).
\end{equation}
\end{lemma}

\begin{proof}
We have $mq=s\,v$ with $v=p_1\cdots p_{j-1}q$ squarefree, $47$-rough and $\omega(v)=j\le J(s)$. So $mq$ is refactorable by the argument in the proof of Proposition~\ref{prop:L}. From $k=mq$ we recover $s=s(k)$ and $v=r(k)$; then $q$ is the largest prime factor of $v$, and $p_1,\dots,p_{j-1}$ are the others. For fixed $m$, the number of admissible primes $q\le x/m$ is
\[
\max\{0,\ \pi(x/m)-\pi(p(m))\}\ \ge\ \pi(x/m)-\pi(p(m)).
\]
\end{proof}

\section{Monotone checkpoints and certified chains}\label{sec:chains}

\begin{lemma}[Monotone checkpoint principle]\label{lem:checkpoint}
Let $1\le a\le b$ be real numbers with $\pi(b)\le2L(a)$. Then $\pi(n)\le2T(n)$ for every real $n\in[a,b]$. If $\pi(b)<2L(a)$, the inequality is strict on $[a,b]$.
\end{lemma}

\begin{proof}
For $n\in[a,b]$, monotonicity and Proposition~\ref{prop:L} give $\pi(n)\le\pi(b)\le2L(a)\le2T(a)\le2T(n)$.
\end{proof}

\begin{definition}
A \emph{certified chain} from $a_0$ to $a_N$ is a sequence of integers $a_0<a_1<\dots<a_N$ with $\pi(a_{i+1})<2L(a_i)$ for $0\le i<N$. By Lemma~\ref{lem:checkpoint}, it proves $\pi(n)<2T(n)$ for all real $n\in[a_0,a_N]$.
\end{definition}

A link certifies an interval of any length, at the cost of one evaluation of $L$ and one of $\pi$. The ratio $2L(x)/\pi(x)$ grows slowly, from $1.135$ at $x=10^6$ to $1.49$ at $x=10^{14}$, so each link can extend the certified range by a factor up to about $2L(a)/\pi(a)$.

We build chains greedily. Given $a_i$, we set $\theta=1.3$ and try $a_{i+1}=\lfloor\theta a_i\rfloor$. As long as $\pi(a_{i+1})>2L(a_i)$, we replace $\theta$ by $1+0.8(\theta-1)$ and try again. For a target endpoint $Y$, the last link is truncated at $\lfloor1.0001\,Y\rfloor+1$.

\begin{proposition}\label{prop:chains}
The sequences in Tables~\ref{tab:chain1} and~\ref{tab:chain2} (Appendix~\ref{app:chains}) are certified chains:
\begin{itemize}
\item $\mathsf C_1$ from $10^{6}$ to $100\,010\,000\,001$, with $58$ links;
\item $\mathsf C_2$ from $10^{11}$ to $100\,010\,000\,000\,001$, with $27$ links.
\end{itemize}
Consequently $\pi(n)<2T(n)$ for every real $n\in[10^{6},\ 1.0001\cdot10^{14}]$.
\end{proposition}

The smallest absolute margin $2L(a_{i-1})-\pi(a_i)$ in $\mathsf C_1$ is $54$ (row $4$ of Table~\ref{tab:chain1}). In $\mathsf C_2$ the relative margin $\bigl(2L(a_{i-1})-\pi(a_i)\bigr)/\pi(a_i)$ is at least $0.0656$ for every link.

\section{The tail \texorpdfstring{$x\ge10^{14}$}{x >= 1e14}}\label{sec:tail}

Fix $X=10^{14}$ and $A=X/\log X=3.102103\ldots\times10^{12}$. Let $\cT_X$ be the set of templates $m$ with $j\le4$ that satisfy
\begin{equation}\label{eq:rule}
\frac{A}{m}>\pi(p(m)).
\end{equation}
This selection rule is a choice. Any finite set of templates with $X/m\ge17$ would be admissible in the argument below. The rule keeps a template exactly when, at $x=X$, its main term $A/m$ exceeds its error term $\pi(p(m))$. Put
\[
W=\sum_{m\in\cT_X}\frac1m,\qquad E=\sum_{m\in\cT_X}\pi(p(m)).
\]
Table~\ref{tab:tail} lists the contributions by the number $j$ of rough prime factors of $mq$.

\begin{table}[ht]
\centering
\caption{The template family $\cT_X$ for $X=10^{14}$.}\label{tab:tail}
\begin{tabular}{crrr}
\toprule
$j$ & number of templates & $W_j=\sum 1/m$ & $E_j=\sum\pi(p(m))$\\
\midrule
1 & 280\,236 & 0.382026077834 & 4\,203\,540\\
2 & 4\,886\,885 & 0.135723931453 & 16\,394\,494\,868\\
3 & 12\,213\,645 & 0.021055159405 & 7\,227\,443\,076\\
4 & 2\,784\,611 & 0.000539288928 & 402\,044\,629\\
\midrule
total & 20\,165\,377 & 0.539344457620 & 24\,028\,186\,113\\
\bottomrule
\end{tabular}
\end{table}

\begin{theorem}[Template tail]\label{thm:tail}
For every real $x\ge10^{14}$,
\[
T(x)>W\,\frac{x}{\log x}-E,\qquad W=0.539344457\ldots,\quad E=24\,028\,186\,113,
\]
and consequently $\pi(x)<2T(x)$.
\end{theorem}

\begin{proof}
Let $x\ge X$ and $m\in\cT_X$. By~\eqref{eq:rule}, $x/m\ge X/m>\pi(p(m))\log X\ge15\log X>17$. Lemma~\ref{lem:RS}(a) gives
\[
\pi(x/m)>\frac{x/m}{\log(x/m)}\ge\frac{x}{m\log x},
\]
and summing~\eqref{eq:template} over $\cT_X$ yields $T(x)>Wx/\log x-E$. By Lemma~\ref{lem:RS}(b) and $\log x\ge\log X$,
\[
\pi(x)<\frac{x}{\log x}\Bigl(1+\frac{3}{2\log X}\Bigr).
\]
Therefore
\[
2T(x)-\pi(x)>\frac{x}{\log x}\,D-2E,\qquad D:=2W-1-\frac{3}{2\log X}=0.0321573636\ldots>0.
\]
Since $x/\log x$ is increasing for $x\ge e$, the right-hand side is at least $A\,D-2E=5.1699\cdot10^{10}>0$.
\end{proof}

\begin{corollary}
$T(x)>0.53\,x/\log x$ for every real $x\ge10^{14}$.
\end{corollary}

\begin{proof}
By Theorem~\ref{thm:tail}, $T(x)-0.53\,x/\log x>(W-0.53)\,x/\log x-E\ge(W-0.53)A-E>2.89\cdot10^{10}-2.41\cdot10^{10}>0$.
\end{proof}

\paragraph{Numerical rigor.} The integer $E$ is computed exactly. The number $W$ is a sum of 20\,165\,377 positive terms. It was computed by two independent programs (Appendices~\ref{code:tail} and~\ref{code:tailpy}):
\begin{itemize}
\item one in $64$-bit-mantissa extended precision with sequential summation;
\item one in IEEE double precision with prefix sums over the primes and a different loop structure.
\end{itemize}
The two values agree to all twelve printed decimals. An a priori bound for the rounding error of either computation is below $10^{-9}$ in absolute value. The slack in the final inequality is $D-2E/A=0.01667\ldots$, more than seven orders of magnitude larger.

The selection rule~\eqref{eq:rule} is evaluated in floating point, so a borderline template might be misclassified. This cannot affect validity: the proof uses only $X/m>17$, and every selected template satisfies $X/m>\pi(p(m))\log X>480$.

\begin{remark}[Why templates with several rough primes are needed]
The templates with $j=1$ contribute only $2W_1=0.7641$, and no family of single-prime templates can do much better. The partial sum of $\sum_{s\in\cS,\,J(s)\ge1}1/s$ over $s\le10^{16}$ is $0.3820268$. By Rankin's trick, the remainder is at most $10^{-8}\prod_{p\le47}(1-p^{-1/2})^{-1}<3\cdot10^{-6}$. So every family of templates with $j=1$ has $2W<0.7641$, whereas the argument of Theorem~\ref{thm:tail} needs $2W>1$. Templates with two and three rough primes supply the missing density. Their total weight grows only slowly with $X$ (roughly like $\log\log X$), which is why the argument starts to work only near $10^{14}$. For instance, at $X=10^{12}$ the same construction gives $2W=1.0077$, which is less than $1+3/(2\log X)=1.0543$.
\end{remark}

\section{\texorpdfstring{Proof of Theorem~\ref{thm:main}}{Proof of the main theorem}}\label{sec:proof}

The interval $[1,\infty)$ is covered by four overlapping pieces.

\begin{enumerate}
\item[(i)] $1\le n\le10^{11}$: exhaustive computation (Section~\ref{sec:exhaustive}). Over this range the minimum of $2T(n)-\pi(n)$ is $0$, and it is first attained at $n=7$. An independent sequential program shows that, for $1\le n\le10^{8}$, equality $\pi(n)=2T(n)$ holds for exactly one $n$, namely $n=7$.
\item[(ii)] $10^{6}\le n\le 1.0001\cdot10^{11}$: chain $\mathsf C_1$ (Proposition~\ref{prop:chains}); the inequality is strict.
\item[(iii)] $10^{11}\le n\le 1.0001\cdot10^{14}$: chain $\mathsf C_2$ (Proposition~\ref{prop:chains}); strict.
\item[(iv)] $n\ge10^{14}$: Theorem~\ref{thm:tail}; strict.
\end{enumerate}

Together these give $\pi(n)\le 2T(n)$ for all $n\ge1$. By (ii)--(iv) the inequality is strict for $n\ge10^6$, and by (i) the only equality below $10^{8}$ occurs at $n=7$. This proves Theorem~\ref{thm:main}. Piece (ii) is not needed for the inequality itself, which already follows from (i), (iii) and (iv). It is used only for the equality statement between $10^8$ and $10^{11}$, and it also gives a second, independent proof of the range $[10^6,10^{11}]$. \qed

\section{Computations}\label{sec:computations}

All computations ran on a shared virtual machine with four cores of an Intel Xeon processor at 2.10\,GHz and 15\,GB of memory. The software was GCC~13.3, Python~3.13, primecount~7.10 and primesieve~12.0~\cite{primecount}. The code is listed in Appendix~\ref{app:code}.

\subsection{Exhaustive verification up to \texorpdfstring{$10^{11}$}{1e11}}\label{sec:exhaustive}

Program \texttt{psieve.c} (Appendix~\ref{code:psieve}) processes $[1,N]$ in chunks in parallel. In each chunk it computes $\tau(k)$ and primality for every $k$ by a segmented sieve: for each prime $p\le\sqrt N$ it divides $p$ out of all multiples of $p$ in the chunk and records the exponent. It records, per chunk:
\begin{itemize}
\item the increments of $T$ and $\pi$;
\item the minimum of the local partial sums of $2\cdot\mathbf 1[\tau(k)\mid k]-\mathbf 1[k \text{ prime}]$.
\end{itemize}
Combining the chunks gives the exact minimum of $2T(n)-\pi(n)$ over $1\le n\le N$. For $N=10^{11}$ this minimum is $0$, first attained at $n=7$; the run took about $41$ minutes. The intermediate values $T(10^{k})$ and $\pi(10^{k})$ ($k\le11$) agree with A111291 and with the classical values of $\pi(10^k)$.

The sequential program \texttt{sieve.c} (Appendix~\ref{code:sieve}) checks $\pi(n)\le2T(n)$ for every $n\le10^{8}$ and counts the equality cases. There is exactly one, $n=7$. Over $100\le n\le10^{8}$ the minimum of $2T(n)/\pi(n)$ is $88/85$, attained at $n=439$. A third program, \texttt{check2.c} (Appendix~\ref{code:check2}), computes $\tau$ by a linear sieve and reproduces $T(10^{8})=3\,558\,181$ and $\pi(10^8)=5\,761\,455$.

\subsection{Evaluation of \texorpdfstring{$L$}{L} and of the chains}

Program \texttt{checkpoints.cpp} (Appendix~\ref{code:checkpoints}) builds a table of $\pi$ up to $Z_0$. The table is a bit array over odd numbers with cumulative counts per $64$-bit word; $Z_0=2\cdot10^{10}$ for $\mathsf C_2$ and $2\cdot10^{9}$ for $\mathsf C_1$. Larger arguments are passed to the primecount library~\cite{primecount}, which by default uses Gourdon's algorithm~\cite{Go01}.

The program enumerates $\cS$ by depth-first search over the exponent vectors of the $15$ primes $\le47$. It evaluates $L(x)$ by Lemma~\ref{lem:phirec}, using the fact that $\Phi_{j+1}(y)=0$ whenever $\Phi_j(y)=0$. It then runs the greedy chain construction of Section~\ref{sec:chains}. Every comparison $\pi(a_{i+1})\le2L(a_i)$ is exact integer arithmetic.

\subsection{Exact values of \texorpdfstring{$T$}{T}}\label{sec:exactdata}

Program \texttt{exactT\_B.cpp} (Appendix~\ref{code:exact}) implements Theorem~\ref{thm:exact} for a general smoothness bound $B$. With $P$ the least prime exceeding $B$, Lemma~\ref{lem:threshold} makes the formula valid for $x<2^{P-1}$; the program refuses larger $x$. Different values of $B$ lead to different decompositions $k=s\,u\,v$, with different sets $\cS$ and $\cU$ and different almost-prime counts. Agreement across several $B$ is therefore a strong consistency check.

The program reproduces every known term $a(0),\dots,a(11)$ of A111291 ($B=47$). It also gives the following values, which agree for every bound $B$ that applies:
\begin{itemize}
\item $T(10^{12})$ for $B=37,43,47,53$;
\item $T(10^{13})$ for $B=43,47,53$;
\item $T(10^{14})$ for $B=47,53$.
\end{itemize}

\begin{table}[ht]
\centering
\caption{$T(10^k)$, $\pi(10^k)$ and the ratio $2T/\pi$. Entries for $k\le11$ agree with OEIS A111291; those for $k=12,13,14$ are new.}\label{tab:values}
\begin{tabular}{rrrr}
\toprule
$k$ & $T(10^k)$ & $\pi(10^k)$ & $2T(10^k)/\pi(10^k)$\\
\midrule
1 & $4$ & $4$ & 2.000000\\
2 & $16$ & $25$ & 1.280000\\
3 & $92$ & $168$ & 1.095238\\
4 & $665$ & $1\,229$ & 1.082181\\
5 & $5\,257$ & $9\,592$ & 1.096122\\
6 & $44\,705$ & $78\,498$ & 1.139010\\
7 & $394\,240$ & $664\,579$ & 1.186435\\
8 & $3\,558\,181$ & $5\,761\,455$ & 1.235168\\
9 & $32\,608\,999$ & $50\,847\,534$ & 1.282619\\
10 & $302\,172\,507$ & $455\,052\,511$ & 1.328078\\
11 & $2\,823\,898\,245$ & $4\,118\,054\,813$ & 1.371472\\
12 & $26\,568\,322\,474^{\ast}$ & $37\,607\,912\,018$ & 1.412911\\
13 & $251\,342\,379\,436^{\ast}$ & $346\,065\,536\,839$ & 1.452571\\
14 & $2\,388\,671\,712\,980^{\ast}$ & $3\,204\,941\,750\,802$ & 1.490618\\
\bottomrule
\end{tabular}
\end{table}

\subsection{Cross-checks}\label{sec:crosscheck}

\begin{enumerate}
\item \emph{Prime counts.} All $85$ values $\pi(a_{i+1})$ used by the chains were recomputed with the primecount command-line tool using the Deléglise--Rivat algorithm (\texttt{--deleglise-rivat}). This is a different algorithm from the library default, and all $85$ values agree.
\item \emph{The lower bound $L$.} Program \texttt{bruteL.c} (Appendix~\ref{code:bruteL}) factors every $k\le N$ and tests the defining condition of $L$ directly. It gives $L(10^6)=44\,559$ and $L(10^8)=3\,553\,787$, equal to the values computed through Lemma~\ref{lem:phirec}.
\item \emph{Exact counts.} Besides the agreement with A111291 for $k\le11$ and with the exhaustive sieve for $k\le11$, \texttt{exactT\_B} agrees with the sieve at the non-round points $4,8,12,16,20\times10^{7}$, and across the smoothness bounds listed in Section~\ref{sec:exactdata}. Consistency $L(x)\le T(x)$ holds at $x=10^{12}$.
\item \emph{Tail certificate.} The C++ program \texttt{tail.cpp} and the independent Python program \texttt{tail\_check.py} agree on every $W_j$ to twelve decimals, and on every template count and $E_j$ exactly. At $X=10^{12}$ both report failure with identical numbers.
\item \emph{Reproducibility.} Both chains were recomputed from scratch, and the outputs are identical to those listed in Appendix~\ref{app:chains}.
\end{enumerate}

\subsection{Timings}\label{sec:timing}

\begin{table}[ht]
\centering
\caption{Wall-clock times on the machine described above (shared with other jobs).}\label{tab:timing}
\begin{tabular}{lll}
\toprule
task & program & time\\
\midrule
exhaustive check, $n\le10^{8}$ (sequential) & \texttt{sieve.c} & 7.5\,s\\
$\tau$ by linear sieve, $n\le10^{8}$ & \texttt{check2.c} & 4.3\,s\\
exhaustive check, $n\le10^{9}$ (4 threads) & \texttt{psieve.c} & 16\,s\\
exhaustive check, $n\le10^{11}$ (4 threads) & \texttt{psieve.c} & about 41\,min\\
chain $\mathsf C_1$ ($10^6\to1.0001\cdot10^{11}$) & \texttt{checkpoints.cpp} & 12\,s\\
chain $\mathsf C_2$ ($10^{11}\to1.0001\cdot10^{14}$) & \texttt{checkpoints.cpp} & 157\,s\\
$T(10^{k})$, $1\le k\le 11$ & \texttt{exactT\_B.cpp} & 11\,s\\
$T(10^{12}),T(10^{13}),T(10^{14})$ & \texttt{exactT\_B.cpp} & 153\,s\\
tail certificate (C++) & \texttt{tail.cpp} & 1.2\,s\\
tail certificate (Python) & \texttt{tail\_check.py} & 5.7\,s\\
\bottomrule
\end{tabular}
\end{table}

The contrast between the exhaustive sieve and the decomposition is the point of the method. Sieving to $10^{11}$ takes about $41$ minutes on four cores. The exact formula gives $T(10^{14})$, a value $1000$ times further out, in under three minutes, and the whole certificate for $[10^{11},10^{14}]$ takes a few minutes.

\section{Discussion}\label{sec:discussion}

\paragraph{What is new.}
To our knowledge the following four ingredients are new:
\begin{itemize}
\item the $\tau$-smoothness threshold (Lemma~\ref{lem:threshold}) used as a reduction device;
\item the exact prime-count formula~\eqref{eq:exact} for $T(x)$;
\item the use of the universal lower bound $L$ inside monotone certified chains;
\item the template tail with an explicit finite certificate.
\end{itemize}
Together they make Colton's conjecture a finite, checkable statement.

\paragraph{Generality.}
The same reasoning applies whenever a condition $f(k)\mid k$ involves a function $f$ that depends only on the exponent pattern of $k$. Examples are $f=\tau^a$, $f(k)=2^{\omega(k)}$, $f(k)=\prod_{p\mid k}(v_p(k)+c)$, and the higher divisor functions $\tau_m$. Below an explicit threshold, $f(k)$ is smooth with respect to a bound depending only on $\log k$. The counting function of $\{k:\ f(k)\mid k\}$ then reduces to prime counts, and certified lower bounds and chains apply verbatim. We expect this to settle other explicit inequalities about such sets with modest computation.

\paragraph{Trust base.}
The proof relies on:
\begin{itemize}
\item the two classical inequalities of Lemma~\ref{lem:RS};
\item the correctness of the compiled programs and of the primecount library.
\end{itemize}
The cross-checks of Section~\ref{sec:crosscheck} are designed so that an error in any single component would very likely have been detected. In particular, two prime-counting algorithms agree on every $\pi$-value used by the chains. The certificates are small: two chains of $85$ links in total, and one inequality with an integer $E$ and a sum $W$. A formal verification in Lean would need formal versions of Lemma~\ref{lem:RS} together with either a verified prime-counting procedure or the exported certificates; we leave this for future work.

\paragraph{Relation to Zelinsky's bound.}
Our tail threshold $10^{14}$ is of the same order as Zelinsky's $7.42\cdot10^{13}$~\cite{Ze02}. We could not consult the full text of~\cite{Ze02} (only its abstract and the summary in~\cite{FC}). Since our argument does not use it, Theorem~\ref{thm:main} is independent of the details of that paper. The computations of Sections~\ref{sec:chains}--\ref{sec:tail} confirm, independently, that no counterexample exists above $7.42\cdot10^{13}$.

\appendix

\section{The certified chains}\label{app:chains}

Row $i$ describes the link from $a_{i-1}$ to $a_i$. It lists the left end $a_{i-1}$, the certified lower bound $L(a_{i-1})\le T(a_{i-1})$, the right end $a_i$, the prime count $\pi(a_i)$, and the margin $2L(a_{i-1})-\pi(a_i)>0$. By Lemma~\ref{lem:checkpoint}, row $i$ proves $\pi(n)<2T(n)$ for all $n\in[a_{i-1},a_i]$.

{\scriptsize
\setlength{\tabcolsep}{4pt}
\begin{longtable}{rrrrrr}
\caption{Chain $\mathsf C_1$ from $10^6$ to $100\,010\,000\,001$ ($58$ links).}\label{tab:chain1}\\
\toprule
$i$ & $a_{i-1}$ & $L(a_{i-1})$ & $a_{i}$ & $\pi(a_{i})$ & $2L(a_{i-1})-\pi(a_{i})$\\
\midrule
\endfirsthead
\toprule
$i$ & $a_{i-1}$ & $L(a_{i-1})$ & $a_{i}$ & $\pi(a_{i})$ & $2L(a_{i-1})-\pi(a_{i})$\\
\midrule
\endhead
1 & 1000000 & 44559 & 1122879 & 87370 & 1748 \\
2 & 1122879 & 49693 & 1260858 & 97245 & 2141 \\
3 & 1260858 & 55414 & 1415792 & 108202 & 2626 \\
4 & 1415792 & 61770 & 1633257 & 123486 & 54 \\
5 & 1633257 & 70678 & 1884125 & 140929 & 427 \\
6 & 1884125 & 80888 & 2173526 & 160879 & 897 \\
7 & 2173526 & 92540 & 2507379 & 183575 & 1505 \\
8 & 2507379 & 105925 & 2892512 & 209610 & 2240 \\
9 & 2892512 & 121267 & 3336801 & 239337 & 3197 \\
10 & 3336801 & 138793 & 3849333 & 273277 & 4309 \\
11 & 3849333 & 158938 & 4440590 & 312035 & 5841 \\
12 & 4440590 & 181997 & 5122664 & 356422 & 7572 \\
13 & 5122664 & 208350 & 5909505 & 407043 & 9657 \\
14 & 5909505 & 238650 & 6817204 & 464971 & 12329 \\
15 & 6817204 & 273328 & 7864326 & 531252 & 15404 \\
16 & 7864326 & 313053 & 9374276 & 625813 & 293 \\
17 & 9374276 & 370005 & 11174136 & 737269 & 2741 \\
18 & 11174136 & 437346 & 13319570 & 868703 & 5989 \\
19 & 13319570 & 516916 & 15876927 & 1023714 & 10118 \\
20 & 15876927 & 611168 & 18925296 & 1206545 & 15791 \\
21 & 18925296 & 722563 & 22558952 & 1422165 & 22961 \\
22 & 22558952 & 854507 & 26890270 & 1676623 & 32391 \\
23 & 26890270 & 1010556 & 32053201 & 1976861 & 44251 \\
24 & 32053201 & 1195187 & 38207415 & 2330927 & 59447 \\
25 & 38207415 & 1413832 & 45543238 & 2749168 & 78496 \\
26 & 45543238 & 1672411 & 54287539 & 3242277 & 102545 \\
27 & 54287539 & 1978644 & 64710746 & 3824526 & 132762 \\
28 & 64710746 & 2341370 & 80241325 & 4682610 & 130 \\
29 & 80241325 & 2877255 & 99499242 & 5734218 & 20292 \\
30 & 99499242 & 3536779 & 123379060 & 7023033 & 50525 \\
31 & 123379060 & 4347348 & 152990034 & 8602884 & 91812 \\
32 & 152990034 & 5345211 & 189707642 & 10540048 & 150374 \\
33 & 189707642 & 6572032 & 235237476 & 12914350 & 229714 \\
34 & 235237476 & 8082080 & 291694470 & 15826422 & 337738 \\
35 & 291694470 & 9939702 & 361701142 & 19397498 & 481906 \\
36 & 361701142 & 12226220 & 448509416 & 23778210 & 674230 \\
37 & 448509416 & 15039304 & 556151675 & 29151098 & 927510 \\
38 & 556151675 & 18503095 & 689628076 & 35743400 & 1262790 \\
39 & 689628076 & 22765533 & 855138814 & 43831785 & 1699281 \\
40 & 855138814 & 28013539 & 1060372129 & 53756452 & 2270626 \\
41 & 1060372129 & 34474062 & 1314861439 & 65936822 & 3011302 \\
42 & 1314861439 & 42429028 & 1709319870 & 84601781 & 256275 \\
43 & 1709319870 & 54661887 & 2222115831 & 108567743 & 756031 \\
44 & 2222115831 & 70433180 & 2888750580 & 139347766 & 1518594 \\
45 & 2888750580 & 90763668 & 3755375754 & 178882751 & 2644585 \\
46 & 3755375754 & 116978945 & 4881988480 & 229668180 & 4289710 \\
47 & 4881988480 & 150783241 & 6346585024 & 294919873 & 6646609 \\
48 & 6346585024 & 194377235 & 8250560531 & 378765377 & 9989093 \\
49 & 8250560531 & 250608026 & 10725728690 & 486519277 & 14696775 \\
50 & 10725728690 & 323137367 & 13943447297 & 625019759 & 21254975 \\
51 & 13943447297 & 416703453 & 18126481486 & 803057904 & 30349002 \\
52 & 18126481486 & 537420518 & 23564425931 & 1031954908 & 42886128 \\
53 & 23564425931 & 693175813 & 30633753710 & 1326266604 & 60085022 \\
54 & 30633753710 & 894165052 & 39823879823 & 1704739831 & 83590273 \\
55 & 39823879823 & 1153536085 & 51771043769 & 2191496567 & 115575603 \\
56 & 51771043769 & 1488288985 & 67302356899 & 2817581748 & 158996222 \\
57 & 67302356899 & 1920366218 & 87493063968 & 3622985305 & 217747131 \\
58 & 87493063968 & 2478098104 & 100010000001 & 4118449214 & 837746994 \\
\bottomrule
\end{longtable}

\begin{longtable}{rrrrrr}
\caption{Chain $\mathsf C_2$ from $10^{11}$ to $100\,010\,000\,000\,001$ ($27$ links).}\label{tab:chain2}\\
\toprule
$i$ & $a_{i-1}$ & $L(a_{i-1})$ & $a_{i}$ & $\pi(a_{i})$ & $2L(a_{i-1})-\pi(a_{i})$\\
\midrule
\endfirsthead
\toprule
$i$ & $a_{i-1}$ & $L(a_{i-1})$ & $a_{i}$ & $\pi(a_{i})$ & $2L(a_{i-1})-\pi(a_{i})$\\
\midrule
\endhead
1 & 100000000000 & 2821796492 & 130000000000 & 5296138250 & 347454734 \\
2 & 130000000000 & 3641815796 & 169000000000 & 6812015990 & 471615602 \\
3 & 169000000000 & 4700523176 & 219700000000 & 8762780454 & 638265898 \\
4 & 219700000000 & 6067511562 & 285610000000 & 11273421612 & 861601512 \\
5 & 285610000000 & 7832681227 & 371293000000 & 14504962265 & 1160400189 \\
6 & 371293000000 & 10112155956 & 482680900000 & 18664796707 & 1559515205 \\
7 & 482680900000 & 13056023905 & 627485170000 & 24020108257 & 2091939553 \\
8 & 627485170000 & 16858178972 & 815730721000 & 30915134830 & 2801223114 \\
9 & 815730721000 & 21769209101 & 1060449937300 & 39793316656 & 3745101546 \\
10 & 1060449937300 & 28112906109 & 1378584918490 & 51226167825 & 4999644393 \\
11 & 1378584918490 & 36307746289 & 1792160394037 & 65950014530 & 6665478048 \\
12 & 1792160394037 & 46894640770 & 2329808512248 & 84913875301 & 8875406239 \\
13 & 2329808512248 & 60572588584 & 3028751065922 & 109340886366 & 11804290802 \\
14 & 3028751065922 & 78245226849 & 3937376385698 & 140807456610 & 15682997088 \\
15 & 3937376385698 & 101080639246 & 5118589301407 & 181345707262 & 20815571230 \\
16 & 5118589301407 & 130588628609 & 6654166091829 & 233575103883 & 27602153335 \\
17 & 6654166091829 & 168721318806 & 8650415919377 & 300872720222 & 36569917390 \\
18 & 8650415919377 & 218002205107 & 11245540695190 & 387592707670 & 48411702544 \\
19 & 11245540695190 & 281694244063 & 14619202903747 & 499348744307 & 64039743819 \\
20 & 14619202903747 & 364016019296 & 19004963774871 & 643379961539 & 84652077053 \\
21 & 19004963774871 & 470422282261 & 24706452907332 & 829020898377 & 111823666145 \\
22 & 24706452907332 & 607966834055 & 32118388779531 & 1068310022059 & 147623646051 \\
23 & 32118388779531 & 785770723938 & 41753905413390 & 1376773153539 & 194768294337 \\
24 & 41753905413390 & 1015629792485 & 54280077037407 & 1774435044640 & 256824540330 \\
25 & 54280077037407 & 1312798803893 & 70564100148629 & 2287125394765 & 338472213021 \\
26 & 70564100148629 & 1697006858491 & 91733330193217 & 2948162225971 & 445851491011 \\
27 & 91733330193217 & 2193771404114 & 100010000000001 & 3205251958942 & 1182290849286 \\
\bottomrule
\end{longtable}
}

\section{Source code}\label{app:code}

The programs below are exactly those that produced the results in this paper.

The programs can be compiled as follows.
{\footnotesize
\begin{verbatim}
gcc -O3 -march=native -fopenmp -o psieve psieve.c -lm
gcc -O2 -o sieve sieve.c -lm
gcc -O2 -o check2 check2.c
gcc -O2 -o bruteL bruteL.c
g++ -O3 -march=native -o checkpoints checkpoints.cpp -lprimecount -lprimesieve
g++ -O3 -march=native -o exactT_B exactT_B.cpp -lprimecount -lprimesieve
g++ -O2 -march=native -o tail tail.cpp
\end{verbatim}
}
The results in this paper were produced by the following invocations.
{\footnotesize
\begin{verbatim}
# exhaustive verification, n <= 1e11 (args: N, chunk size, output table)
./psieve 100000000000 50000000 chunks.txt
# sequential verification with equality count, n <= 1e8
./sieve 100000000
# chains (args: Z0, prime-list limit, start, end, initial step ratio)
./checkpoints 2000000000 20000000 1000000 100000000000 1.3              # C1
./checkpoints 20000000000 20000000 100000000000 100000000000000 1.3     # C2
# tail certificate (args: X, prime-list limit, max j)
./tail 1e14 50000000 4
python3 tail_check.py 1e14 20000000
# exact values (args: B, Z0, prime-list limit, x_1, x_2, ...)
./exactT_B 47 20000000000 30000000 1000000000000 10000000000000 \
           100000000000000
\end{verbatim}
}

\subsection{Exhaustive parallel verification (\texttt{psieve.c})}\label{code:psieve}
\begin{lstlisting}[language=C]
// Parallel exhaustive verification of pi(n) <= 2*T(n) for all 1 <= n <= N.
// T(n) = #{k<=n : tau(k) | k}. Chunks processed in parallel; per-chunk we record
// dT, dP and the minimum of (2*T_loc - P_loc) over prefixes inside the chunk, then combine.
// Also records T and pi at checkpoints (powers of 10 and multiples of chunk).
#include <stdio.h>
#include <stdlib.h>
#include <stdint.h>
#include <string.h>
#include <math.h>
#include <omp.h>
typedef uint64_t u64; typedef int64_t i64; typedef uint32_t u32;
static u32 *pr; static u64 np;
typedef struct { u64 dT, dP; i64 minloc; u64 argmin; } chunkres;
static void do_chunk(u64 lo, u64 hi, u64 *rem, u32 *tau, chunkres *out) {
    u64 len = hi - lo + 1;
    for (u64 i = 0; i < len; i++) { rem[i] = lo + i; tau[i] = 1; }
    for (u64 t = 0; t < np; t++) {
        u64 p = pr[t]; if (p * p > hi) break;
        u64 start = ((lo + p - 1) / p) * p;
        for (u64 m = start; m <= hi; m += p) {
            u64 i = m - lo; u32 e = 0; u64 r = rem[i];
            do { r /= p; e++; } while (r % p == 0);
            rem[i] = r; tau[i] *= (e + 1);
        }
    }
    u64 T = 0, P = 0; i64 minloc = INT64_MAX; u64 argmin = lo;
    for (u64 i = 0; i < len; i++) {
        u64 k = lo + i; u32 tk = tau[i]; if (rem[i] > 1) tk *= 2;
        if (k >= 2 && tk == 2) P++;
        if (k % tk == 0) T++;
        i64 v = 2 * (i64)T - (i64)P;
        if (v < minloc) { minloc = v; argmin = k; }
    }
    out->dT = T; out->dP = P; out->minloc = minloc; out->argmin = argmin;
}
int main(int argc, char **argv) {
    u64 N = strtoull(argv[1], 0, 10);
    u64 CH = argc > 2 ? strtoull(argv[2], 0, 10) : 100000000ULL;
    u64 R = (u64)sqrtl((long double)N) + 2;
    char *isc = calloc(R + 1, 1); pr = malloc(sizeof(u32) * (R + 1)); np = 0;
    for (u64 i = 2; i <= R; i++) { if (!isc[i]) { pr[np++] = (u32)i; for (u64 j = i * i; j <= R; j += i) isc[j] = 1; } }
    u64 nch = (N + CH - 1) / CH;
    chunkres *res = calloc(nch, sizeof(chunkres));
    #pragma omp parallel
    {
        u64 *rem = malloc(sizeof(u64) * CH); u32 *tau = malloc(sizeof(u32) * CH);
        #pragma omp for schedule(dynamic, 1)
        for (u64 c = 0; c < nch; c++) {
            u64 lo = c * CH + 1, hi = lo + CH - 1; if (hi > N) hi = N;
            do_chunk(lo, hi, rem, tau, &res[c]);
        }
        free(rem); free(tau);
    }
    u64 T = 0, P = 0; i64 globalmin = 1LL << 62; u64 gargmin = 0;
    double minratio = 1e18; u64 argr = 0;
    for (u64 c = 0; c < nch; c++) {
        i64 base = 2 * (i64)T - (i64)P;
        i64 m = base + res[c].minloc;
        if (m < globalmin) { globalmin = m; gargmin = res[c].argmin; }
        T += res[c].dT; P += res[c].dP;
        u64 hi = (c + 1) * CH; if (hi > N) hi = N;
        double r = (P > 0) ? 2.0 * (double)T / (double)P : 1e18;
        if (r < minratio && hi >= 1000000) { minratio = r; argr = hi; }
        // print checkpoint at powers of ten
        u64 x = hi; int pow10 = 1; u64 t = x; while (t > 1) { if (t % 10) { pow10 = 0; break; } t /= 10; }
        if (pow10 || c == nch - 1) printf("x=%llu T=%llu pi=%llu 2T-pi=%lld ratio=%.6f\n", (unsigned long long)x, (unsigned long long)T, (unsigned long long)P, (long long)(2 * (i64)T - (i64)P), r);
    }
    printf("RESULT N=%llu: min over 1<=n<=N of (2T(n)-pi(n)) = %lld at n=%llu  %s\n", (unsigned long long)N, (long long)globalmin, (unsigned long long)gargmin, globalmin >= 0 ? "(conjecture holds)" : "(COUNTEREXAMPLE)");
    // dump chunk-end table for later use
    FILE *f = fopen(argc > 3 ? argv[3] : "chunks.txt", "w");
    T = 0; P = 0;
    for (u64 c = 0; c < nch; c++) { T += res[c].dT; P += res[c].dP; u64 hi = (c + 1) * CH; if (hi > N) hi = N; fprintf(f, "%llu %llu %llu\n", (unsigned long long)hi, (unsigned long long)T, (unsigned long long)P); }
    fclose(f);
    return 0;
}
\end{lstlisting}

\subsection{Sequential verification with equality count (\texttt{sieve.c})}\label{code:sieve}
\begin{lstlisting}[language=C]
// Exhaustive verification of Colton's conjecture pi(n) <= 2 T(n) for all n <= N,
// where T(n) = #{k <= n : tau(k) | k}. Segmented sieve computing tau(k) and primality.
// Usage: ./sieve N [segment]
#include <stdio.h>
#include <stdlib.h>
#include <stdint.h>
#include <string.h>
#include <math.h>
typedef uint64_t u64; typedef uint32_t u32;
int main(int argc, char **argv) {
    u64 N = strtoull(argv[1], 0, 10);
    u64 SEG = argc > 2 ? strtoull(argv[2], 0, 10) : (1ULL << 24);
    u64 R = (u64)sqrtl((long double)N) + 2;
    // base primes up to R
    char *isc = calloc(R + 1, 1);
    u32 *pr = malloc(sizeof(u32) * (R + 1)); u64 np = 0;
    for (u64 i = 2; i <= R; i++) { if (!isc[i]) { pr[np++] = (u32)i; for (u64 j = i * i; j <= R; j += i) isc[j] = 1; } }
    u64 *rem = malloc(sizeof(u64) * SEG);   // remaining cofactor
    u32 *tau = malloc(sizeof(u32) * SEG);
    u64 T = 0, P = 0; // counts
    double minratio = 1e9; u64 argmin = 0; u64 tight = 0;
    // record checkpoints
    for (u64 lo = 1; lo <= N; lo += SEG) {
        u64 hi = lo + SEG - 1; if (hi > N) hi = N;
        u64 len = hi - lo + 1;
        for (u64 i = 0; i < len; i++) { rem[i] = lo + i; tau[i] = 1; }
        for (u64 t = 0; t < np; t++) {
            u64 p = pr[t]; if (p * p > hi) break;
            u64 start = ((lo + p - 1) / p) * p;
            for (u64 m = start; m <= hi; m += p) {
                u64 i = m - lo; u32 e = 0;
                while (rem[i] % p == 0) { rem[i] /= p; e++; }
                tau[i] *= (e + 1);
            }
        }
        for (u64 i = 0; i < len; i++) {
            u64 k = lo + i;
            u32 tk = tau[i]; if (rem[i] > 1) tk *= 2;
            if (k >= 2 && tk == 2) P++;
            if (k % tk == 0) T++;
            // check pi(k) <= 2 T(k)
            if (P > 2 * T) { printf("COUNTEREXAMPLE at n=%llu: pi=%llu T=%llu\n", (unsigned long long)k, (unsigned long long)P, (unsigned long long)T); fflush(stdout); }
            if (P > 0) {
                double r = (2.0 * (double)T) / (double)P;
                if (k >= 100 && r < minratio) { minratio = r; argmin = k; }
                if (2 * T == P) tight++;
            }
        }
        // checkpoint output at powers of 10
    }
    printf("N=%llu pi=%llu T=%llu  min 2T/pi (n>=100) = %.6f at n=%llu ; equality cases: %llu\n",
           (unsigned long long)N, (unsigned long long)P, (unsigned long long)T, minratio, (unsigned long long)argmin, (unsigned long long)tight);
    return 0;
}
\end{lstlisting}

\subsection{\texorpdfstring{Linear-sieve cross-check of $T$ and $\pi$ (\texttt{check2.c})}{Linear-sieve cross-check of T and pi (check2.c)}}\label{code:check2}
\begin{lstlisting}[language=C]
// independent check: tau via linear sieve (smallest prime factor) up to N
#include <stdio.h>
#include <stdlib.h>
#include <stdint.h>
int main(int argc,char**argv){
  long N=atol(argv[1]);
  int *spf=calloc(N+1,sizeof(int)); unsigned char *e=calloc(N+1,1); int *tau=calloc(N+1,sizeof(int)); int *rest=calloc(N+1,sizeof(int));
  int *pr=malloc(sizeof(int)*(N/2+10)); long np=0;
  tau[1]=1;
  for(long i=2;i<=N;i++){
    if(!spf[i]){spf[i]=i; pr[np++]=i; e[i]=1; tau[i]=2; rest[i]=1;}
    for(long j=0;j<np;j++){ long p=pr[j]; if(p>spf[i]||p*i>N) break; long m=p*i; spf[m]=p;
      if(p==spf[i]){ e[m]=e[i]+1; rest[m]=rest[i]; tau[m]=tau[rest[i]]*(e[m]+1);} else { e[m]=1; rest[m]=i; tau[m]=tau[i]*2; } }
  }
  long T=0,P=0; for(long k=1;k<=N;k++){ if(k%tau[k]==0) T++; if(tau[k]==2) P++; }
  printf("N=%ld T=%ld pi=%ld\n",N,T,P);
}
\end{lstlisting}

\subsection{\texorpdfstring{Brute-force evaluation of $L$ (\texttt{bruteL.c})}{Brute-force evaluation of L (bruteL.c)}}\label{code:bruteL}
\begin{lstlisting}[language=C]
// Brute-force L(N): count k <= N with k = s*r, s 47-smooth, r squarefree with all primes > 47,
// omega(r) = j and 2^j * tau(s) | s. Independent check of the structured formula for L.
#include <stdio.h>
#include <stdlib.h>
#include <stdint.h>
int main(int argc, char **argv) {
    uint32_t N = (uint32_t)strtoul(argv[1], 0, 10);
    uint32_t *spf = calloc((size_t)N + 1, sizeof(uint32_t));
    for (uint64_t i = 2; i <= N; i++) if (!spf[i]) for (uint64_t k = i; k <= N; k += i) if (!spf[k]) spf[k] = (uint32_t)i;
    uint64_t L = 0, T = 0;
    for (uint32_t k = 1; k <= N; k++) {
        uint32_t m = k; uint64_t s = 1, taus = 1, tauk = 1; int j = 0, sqf = 1;
        while (m > 1) {
            uint32_t p = spf[m]; int e = 0; uint64_t pe = 1;
            while (m % p == 0) { m /= p; e++; pe *= p; }
            tauk *= (e + 1);
            if (p <= 47) { s *= pe; taus *= (e + 1); }
            else { if (e > 1) sqf = 0; j++; }
        }
        if (k % tauk == 0) T++;
        if (sqf && s % ((taus << j)) == 0) L++;
    }
    printf("N=%u bruteL=%llu T=%llu\n", N, (unsigned long long)L, (unsigned long long)T);
    return 0;
}
\end{lstlisting}

\subsection{\texorpdfstring{Lower bound $L$ and certified chains (\texttt{checkpoints.cpp})}{Lower bound L and certified chains (checkpoints.cpp)}}\label{code:checkpoints}
\begin{lstlisting}[language=C++]
// Rigorous lower bound L(x) <= T(x) = #{k <= x : tau(k) | k}.
// k = s*r, s 47-smooth with tau(s) | s, r squarefree with j prime factors all > 47,
// counted only when 2^j * tau(s) | s  (then tau(k) = tau(s)*2^j divides s, hence k).
// Distinct (s, r) give distinct k by unique factorization.
// L(x) = sum_{s} sum_{j=0}^{J(s)} Phi_j(floor(x/s)),  Phi_j(y) = #{r<=y squarefree, omega(r)=j, all p>47}.
// pi(z): exact via bit table for z <= Z0, else primecount::pi.
#include <bits/stdc++.h>
#include <primecount.hpp>
using namespace std;
typedef long long ll; typedef unsigned long long ull;

static ull Z0;                   // table limit
static vector<ull> bits;         // odd-only bitset: bit i <-> number 2i+1 is prime
static vector<uint32_t> blockcnt; // prefix count of primes (odd) per 64-bit word
static vector<uint32_t> smallp;  // primes up to PL for iteration
static unordered_map<ull, ll> pcache;

void build_table(ull Z) {
    Z0 = Z;
    ull nbits = Z / 2 + 1; ull nw = nbits / 64 + 2;
    bits.assign(nw, ~0ULL);
    // bit i: number 2i+1. mark 1 (i=0) composite
    auto clr = [&](ull i) { bits[i >> 6] &= ~(1ULL << (i & 63)); };
    clr(0);
    ull r = (ull)sqrtl((long double)Z) + 1;
    vector<char> sm(r + 1, 1);
    for (ull p = 3; p <= r; p += 2) if (sm[p]) {
        for (ull q = p * p; q <= r; q += 2 * p) sm[q] = 0;
        for (ull m = p * p; m <= Z; m += 2 * p) clr(m / 2);
    }
    // clear bits beyond Z
    for (ull i = nbits; i < nw * 64; i++) clr(i);
    blockcnt.assign(nw + 1, 0);
    ull c = 0;
    for (ull w = 0; w < nw; w++) { blockcnt[w] = (uint32_t)(c & 0xffffffffULL); c += __builtin_popcountll(bits[w]); }
    // store high parts separately: counts up to ~1e9 fit in 32 bits for Z<=2e10? pi(2e10)~8.8e8 < 4.29e9 OK
}
inline ll pi_table(ull z) { // z <= Z0
    if (z < 2) return 0;
    if (z < 3) return 1;
    ull i = (z - 1) / 2; // number 2i+1 <= z, largest odd <= z
    ull w = i >> 6; ull b = i & 63;
    ull mask = (b == 63) ? ~0ULL : ((1ULL << (b + 1)) - 1);
    return 1 + (ll)blockcnt[w] + __builtin_popcountll(bits[w] & mask); // +1 for prime 2
}
ll PI(ull z) {
    if (z <= Z0) return pi_table(z);
    auto it = pcache.find(z); if (it != pcache.end()) return it->second;
    ll v = primecount::pi((int64_t)z);
    pcache[z] = v; return v;
}
static const ll PI47 = 15; // pi(47)
// Phi_j(y, idx): number of squarefree r <= y, omega(r)=j, all prime factors >= smallp[idx] (> 47)
ll Phi(ull y, int j, size_t idx) {
    if (j == 0) return y >= 1 ? 1 : 0;
    if (idx >= smallp.size()) { fprintf(stderr, "prime list too short\n"); exit(1); }
    ull p0 = smallp[idx];
    if (j == 1) { ll v = PI(y) - PI(p0 - 1); return v > 0 ? v : 0; }
    ll total = 0;
    for (size_t t = idx; t < smallp.size(); t++) {
        ull p = smallp[t];
        // need p * (next j-1 primes, each > p) <= y ; lower bound check p^j <= y
        long double pj = powl((long double)p, j);
        if (pj > (long double)y) break;
        total += Phi(y / p, j - 1, t + 1);
    }
    return total;
}

ll computeL(ull x, vector<pair<ull,int>> &good) {
    ll L = 0;
    for (auto &g : good) {
        if (g.first > x) continue;
        ull y = x / g.first;
        for (int j = 0; j <= g.second; j++) {
            ll f = Phi(y, j, 0);
            if (f == 0 && j > 0) break;
            L += f;
        }
    }
    return L;
}
int main(int argc, char **argv) {
    // args: Z0 PLIMIT x0 Xend q
    ull Z = strtoull(argv[1], 0, 10);
    ull PL = strtoull(argv[2], 0, 10);
    ull x0 = strtoull(argv[3], 0, 10);
    ull Xend = strtoull(argv[4], 0, 10);
    double q = atof(argv[5]);
    build_table(Z);
    { vector<char> sv(PL + 1, 1);
      for (ull i = 2; i * i <= PL; i++) if (sv[i]) for (ull k = i * i; k <= PL; k += i) sv[k] = 0;
      for (ull i = 53; i <= PL; i++) if (sv[i]) smallp.push_back((uint32_t)i); }
    const int P[15] = {2,3,5,7,11,13,17,19,23,29,31,37,41,43,47};
    vector<pair<ull,int>> good;
    ull X = (ull)(Xend * 1.5);
    function<void(int, ull, ull)> dfs = [&](int i, ull s, ull tau) {
        if (i == 15) { if (s % tau == 0) { ull qq = s / tau; int J = 0; while (qq % 2 == 0) { qq /= 2; J++; } good.push_back({s, J}); } return; }
        ull v = s; int e = 0;
        while (true) { dfs(i + 1, v, tau * (e + 1)); if (v > X / P[i]) break; v *= P[i]; e++; }
    };
    dfs(0, 1, 1);
    sort(good.begin(), good.end());
    fprintf(stderr, "good s: %zu\n", good.size());
    ull x = x0; ll L = computeL(x, good);
    printf("CHECKPOINT x=%llu L=%lld pi=%lld 2L/pi=%.6f\n", x, L, PI(x), 2.0 * L / PI(x)); fflush(stdout);
    int step = 0;
    while (x < Xend) {
        // find next checkpoint: largest y = x*q' (q' <= q) with pi(y) <= 2L
        double qq = q; ull y; ll py;
        while (true) {
            y = (ull)((long double)x * qq);
            if (y > Xend * 1.0001 && x < Xend) y = (ull)(Xend * 1.0001) + 1;
            py = PI(y);
            if (py <= 2 * L) break;
            qq = 1.0 + (qq - 1.0) * 0.8;
            if (qq < 1.0001) { printf("FAILED to advance at x=%llu\n", x); return 1; }
        }
        printf("  interval [%llu, %llu]: pi(right)=%lld <= 2*L(left)=%lld  OK\n", x, y, py, 2 * L); fflush(stdout);
        x = y; L = computeL(x, good); step++;
        printf("CHECKPOINT x=%llu L=%lld pi=%lld 2L/pi=%.6f (step %d)\n", x, L, PI(x), 2.0 * L / PI(x), step); fflush(stdout);
    }
    printf("VERIFIED: pi(n) <= 2T(n) for all n in [%llu, %llu]\n", x0, x);
    return 0;
}
\end{lstlisting}

\subsection{\texorpdfstring{Exact formula for $T(x)$ with smoothness bound $B$ (\texttt{exactT\_B.cpp})}{Exact formula for T(x) with smoothness bound B (exactT\_B.cpp)}}\label{code:exact}
\begin{lstlisting}[language=C++]
// Exact T(x) = #{k <= x : tau(k) | k} for x < 2^47 via the smooth-rough decomposition.
// k = s * u * v with s the 47-smooth part, u the powerful part of the 47-rough part
// (primes > 47 with exponent >= 2) and v the squarefree part (primes > 47 with exponent 1).
// For k <= x < 2^47 every factor e+1 of tau(k) is <= 47 (exponents of 2..47 are <= 46,
// exponents of primes >= 53 are <= 8), so tau(k) is 47-smooth and coprime to u*v. Hence
//   tau(k) | k  <=>  tau(s) * tau(u) * 2^omega(v) | s.
// T(x) = sum_{s: tau(s)|s} sum_{u: tau(u) | d(s)} sum_{j=0}^{v2(d(s)/tau(u))} Phi_j^{(u)}(x/(s u)),
// d(s) = s / tau(s), Phi_j^{(u)}(y) = #{v <= y squarefree, primes > 47, omega(v) = j, gcd(v,u) = 1}.
#include <bits/stdc++.h>
#include <primecount.hpp>
using namespace std;
typedef long long ll; typedef unsigned long long ull;

static ull Z0;
static vector<ull> bits;
static vector<uint32_t> blockcnt;
static vector<uint32_t> smallp;
static unordered_map<ull, ll> pcache;

void build_table(ull Z) {
    Z0 = Z;
    ull nbits = Z / 2 + 1; ull nw = nbits / 64 + 2;
    bits.assign(nw, ~0ULL);
    auto clr = [&](ull i) { bits[i >> 6] &= ~(1ULL << (i & 63)); };
    clr(0);
    ull r = (ull)sqrtl((long double)Z) + 1;
    vector<char> sm(r + 1, 1);
    for (ull p = 3; p <= r; p += 2) if (sm[p]) {
        for (ull q = p * p; q <= r; q += 2 * p) sm[q] = 0;
        for (ull m = p * p; m <= Z; m += 2 * p) clr(m / 2);
    }
    for (ull i = nbits; i < nw * 64; i++) clr(i);
    blockcnt.assign(nw + 1, 0);
    ull c = 0;
    for (ull w = 0; w < nw; w++) { blockcnt[w] = (uint32_t)c; c += __builtin_popcountll(bits[w]); }
    if (c > 0xffffffffULL) { fprintf(stderr, "table too large for 32-bit counts\n"); exit(1); }
}
inline ll pi_table(ull z) {
    if (z < 2) return 0;
    if (z < 3) return 1;
    ull i = (z - 1) / 2; ull w = i >> 6; ull b = i & 63;
    ull mask = (b == 63) ? ~0ULL : ((1ULL << (b + 1)) - 1);
    return 1 + (ll)blockcnt[w] + __builtin_popcountll(bits[w] & mask);
}
ll PI(ull z) {
    if (z <= Z0) return pi_table(z);
    auto it = pcache.find(z); if (it != pcache.end()) return it->second;
    ll v = primecount::pi((int64_t)z);
    pcache[z] = v; return v;
}
// squarefree r <= y, omega(r) = j, all prime factors >= smallp[idx]
ll Phi(ull y, int j, size_t idx) {
    if (j == 0) return y >= 1 ? 1 : 0;
    if (idx >= smallp.size()) { fprintf(stderr, "prime list too short\n"); exit(1); }
    ull p0 = smallp[idx];
    if (j == 1) { ll v = PI(y) - PI(p0 - 1); return v > 0 ? v : 0; }
    ll total = 0;
    for (size_t t = idx; ; t++) {
        if (t + 1 >= smallp.size()) { fprintf(stderr, "prime list too short\n"); exit(1); }
        ull p = smallp[t];
        long double pj = powl((long double)p, j);
        if (pj > (long double)y) break;
        total += Phi(y / p, j - 1, t + 1);
    }
    return total;
}
// same, additionally coprime to the primes Q[0..nq-1]
ll PhiQ(ull y, int j, const ull *Q, int nq) {
    if (j == 0) return y >= 1 ? 1 : 0;
    if (y == 0) return 0;
    if (nq == 0) return Phi(y, j, 0);
    ull q = Q[nq - 1];
    return PhiQ(y, j, Q, nq - 1) - PhiQ(y / q, j - 1, Q, nq);
}
struct Pow { ull u; ull tau; int nq; ull Q[8]; };

int main(int argc, char **argv) {
    // args: B Z0 PL x1 x2 ...
    int B = atoi(argv[1]);
    ull Z = strtoull(argv[2], 0, 10);
    ull PL = strtoull(argv[3], 0, 10);
    vector<ull> xs; for (int a = 4; a < argc; a++) xs.push_back(strtoull(argv[a], 0, 10));
    ull X = *max_element(xs.begin(), xs.end());
    vector<int> Pv; int Pnext = 0;
    for (int c = 2; ; c++) { bool pr = true; for (int d = 2; d * d <= c; d++) if (c % d == 0) { pr = false; break; } if (pr) { if (c <= B) Pv.push_back(c); else { Pnext = c; break; } } }
    if (Pnext - 1 >= 64 || X >= (1ULL << (Pnext - 1))) { fprintf(stderr, "x must be < 2^%d\n", Pnext - 1); return 1; }
    int NP = (int)Pv.size();
    build_table(Z);
    { vector<char> sv(PL + 1, 1);
      for (ull i = 2; i * i <= PL; i++) if (sv[i]) for (ull k = i * i; k <= PL; k += i) sv[k] = 0;
      for (ull i = (ull)B + 1; i <= PL; i++) if (sv[i]) smallp.push_back((uint32_t)i); }
    // good s: 47-smooth, tau(s) | s
    vector<pair<ull,ull>> good; // (s, d = s/tau(s))
    function<void(int, ull, ull)> dfs = [&](int i, ull s, ull tau) {
        if (i == NP) { if (s % tau == 0) good.push_back({s, s / tau}); return; }
        ull v = s; int e = 0;
        while (true) { dfs(i + 1, v, tau * (e + 1)); if (v > X / Pv[i]) break; v *= Pv[i]; e++; }
    };
    dfs(0, 1, 1);
    sort(good.begin(), good.end());
    // powerful 47-rough u <= X/9 (u > 1 needs tau(u) >= 3 dividing d(s), and d(s) <= 2 for s < 9)
    ull UMAX = X / 9;
    vector<Pow> pw;
    { Pow one; one.u = 1; one.tau = 1; one.nq = 0; pw.push_back(one); }
    function<void(size_t, ull, ull, Pow &)> dfsu = [&](size_t t, ull u, ull tau, Pow &cur) {
        for (size_t i = t; ; i++) {
            if (i >= smallp.size()) { fprintf(stderr, "prime list too short for u\n"); exit(1); }
            ull q = smallp[i];
            if ((long double)u * q * q > (long double)UMAX) break;
            ull v = u * q * q; int f = 2;
            while (true) {
                Pow nx = cur; nx.Q[nx.nq++] = q; nx.u = v; nx.tau = tau * (f + 1);
                pw.push_back(nx);
                dfsu(i + 1, v, tau * (f + 1), nx);
                if ((long double)v * q > (long double)UMAX) break;
                v *= q; f++;
            }
        }
    };
    { Pow root; root.u = 1; root.tau = 1; root.nq = 0; dfsu(0, 1, 1, root); }
    sort(pw.begin(), pw.end(), [](const Pow &a, const Pow &b) { return a.u < b.u; });
    fprintf(stderr, "good s: %zu, powerful rough u: %zu\n", good.size(), pw.size());
    for (ull x : xs) {
        ll T = 0;
        for (auto &g : good) {
            ull s = g.first, d = g.second;
            if (s > x) break;
            ull xs_ = x / s;
            for (auto &w : pw) {
                if (w.u > xs_) break;
                if (d % w.tau) continue;
                ull r = d / w.tau; int J = 0; while (r % 2 == 0) { r /= 2; J++; }
                ull y = xs_ / w.u;
                for (int j = 0; j <= J; j++) {
                    ll f = PhiQ(y, j, w.Q, w.nq);
                    if (f == 0 && j > 0) break;
                    T += f;
                }
            }
        }
        printf("x=%llu T(x)=%lld pi(x)=%lld 2T-pi=%lld 2T/pi=%.6f\n", x, T, PI(x), 2 * T - PI(x), 2.0 * T / PI(x));
        fflush(stdout);
    }
    return 0;
}
\end{lstlisting}

\subsection{Template tail certificate (\texttt{tail.cpp})}\label{code:tail}
\begin{lstlisting}[language=C++]
// Analytic tail certificate for Colton's conjecture.
// For x >= X:  T(x) >= sum_{m in M} ( pi(x/m) - pi(p(m)) )
// where m = s * p_1 ... p_{j-1}, s 47-smooth with 2^j tau(s) | s, 53 <= p_1 < ... < p_{j-1} primes,
// p(m) = p_{j-1} (or 47 if j = 1), and the counted integers are k = m*q with q prime, p(m) < q <= x/m.
// Each such k has tau(k) = 2^j tau(s) | s | k, and k determines (s, p_1..p_{j-1}, q) uniquely.
// Using pi(y) >= y/log y (y >= 17) and log(x/m) <= log x:
//   T(x) >= (x/log x) * W - E,   W = sum 1/m,  E = sum pi(p(m)).
// Using pi(x) < (x/log x)(1 + 3/(2 log x)) (x > 1):
//   2T(x) - pi(x) > (x/log x)(2W - 1 - 3/(2 log X)) - 2E, increasing in x once 2W - 1 - 3/(2 log X) > 0.
// So it suffices that (X/log X)(2W - 1 - 3/(2 log X)) >= 2E.
// m is included iff A/m > pi(p(m)) with A = X/log X (any choice is valid; this one is near-optimal),
// which also guarantees x/m >= X/m > 15 log X >= 17 for every included m.
#include <bits/stdc++.h>
using namespace std;
typedef unsigned long long ull; typedef long double ld;

int main(int argc, char **argv) {
    ld X = strtold(argv[1], 0);
    ull PMAX = argc > 2 ? strtoull(argv[2], 0, 10) : 100000000ULL;
    int JMAX = argc > 3 ? atoi(argv[3]) : 5;
    ld lX = logl(X), A = X / lX;
    // primes >= 53 up to PMAX
    vector<char> sv(PMAX + 1, 1); sv[0] = sv[1] = 0;
    for (ull i = 2; i * i <= PMAX; i++) if (sv[i]) for (ull k = i * i; k <= PMAX; k += i) sv[k] = 0;
    vector<ull> q; for (ull i = 53; i <= PMAX; i++) if (sv[i]) q.push_back(i);
    size_t NQ = q.size();
    vector<ld> inv(NQ + 1, 0); vector<ull> pis(NQ + 1, 0);
    for (size_t t = 0; t < NQ; t++) { inv[t + 1] = inv[t] + 1.0L / (ld)q[t]; pis[t + 1] = pis[t] + (ull)(t + 16); }
    // pi(q[t]) = t + 16 since pi(53) = 16
    auto first_fail = [&](ld B, size_t lo) -> size_t { // smallest t >= lo with B/q[t] <= t+16 (monotone)
        size_t a = lo, b = NQ; // invariant: answer in [a, b]
        while (a < b) { size_t mid = (a + b) / 2; if (B / (ld)q[mid] > (ld)(mid + 16)) a = mid + 1; else b = mid; }
        if (a == NQ) { fprintf(stderr, "PMAX too small\n"); exit(1); }
        return a;
    };
    // enumerate 47-smooth s with tau(s) | s, J(s) = v2(s/tau(s)) >= 1, and A/s > 15
    const int P[15] = {2,3,5,7,11,13,17,19,23,29,31,37,41,43,47};
    ld SMAXL = A / 15; ull SMAX = (ull)SMAXL;
    vector<pair<ull,int>> good;
    function<void(int, ull, ull)> dfs = [&](int i, ull s, ull tau) {
        if (i == 15) { if (s % tau == 0) { ull r = s / tau; int J = 0; while (r % 2 == 0) { r /= 2; J++; } if (J >= 1) good.push_back({s, J}); } return; }
        ull v = s; int e = 0;
        while (true) { dfs(i + 1, v, tau * (e + 1)); if (v > SMAX / P[i]) break; v *= P[i]; e++; }
    };
    dfs(0, 1, 1);
    sort(good.begin(), good.end());
    ld W[8] = {0}; ull E[8] = {0}; ull cnt[8] = {0};
    for (auto &g : good) {
        ull s = g.first; int J = g.second;
        ld Bs = A / (ld)s;
        if (Bs > 15) { W[1] += 1.0L / (ld)s; E[1] += 15; cnt[1]++; }
        if (J >= 2 && JMAX >= 2) {
            size_t t2 = first_fail(Bs, 0);
            W[2] += inv[t2] / (ld)s; E[2] += pis[t2]; cnt[2] += t2;
        }
        if (J >= 3 && JMAX >= 3) {
            for (size_t t1 = 0; t1 < NQ; t1++) {
                ld B1 = Bs / (ld)q[t1];
                size_t c = first_fail(B1, t1 + 1);
                if (c == t1 + 1) break;
                W[3] += (inv[c] - inv[t1 + 1]) / ((ld)s * (ld)q[t1]); E[3] += pis[c] - pis[t1 + 1]; cnt[3] += c - t1 - 1;
            }
        }
        if (J >= 4 && JMAX >= 4) {
            for (size_t t1 = 0; t1 < NQ; t1++) {
                ld B1 = Bs / (ld)q[t1];
                bool any1 = false;
                for (size_t t2 = t1 + 1; t2 < NQ; t2++) {
                    ld B2 = B1 / (ld)q[t2];
                    size_t c = first_fail(B2, t2 + 1);
                    if (c == t2 + 1) break;
                    any1 = true;
                    W[4] += (inv[c] - inv[t2 + 1]) / ((ld)s * (ld)q[t1] * (ld)q[t2]); E[4] += pis[c] - pis[t2 + 1]; cnt[4] += c - t2 - 1;
                }
                if (!any1) break;
            }
        }
        if (J >= 5 && JMAX >= 5) {
            for (size_t t1 = 0; t1 < NQ; t1++) {
                ld B1 = Bs / (ld)q[t1]; bool any1 = false;
                for (size_t t2 = t1 + 1; t2 < NQ; t2++) {
                    ld B2 = B1 / (ld)q[t2]; bool any2 = false;
                    for (size_t t3 = t2 + 1; t3 < NQ; t3++) {
                        ld B3 = B2 / (ld)q[t3];
                        size_t c = first_fail(B3, t3 + 1);
                        if (c == t3 + 1) break;
                        any2 = true;
                        W[5] += (inv[c] - inv[t3 + 1]) / ((ld)s * (ld)q[t1] * (ld)q[t2] * (ld)q[t3]); E[5] += pis[c] - pis[t3 + 1]; cnt[5] += c - t3 - 1;
                    }
                    if (!any2) break; any1 = true;
                }
                if (!any1) break;
            }
        }
    }
    ld Wt = 0; ull Et = 0;
    for (int j = 1; j <= 5; j++) { Wt += W[j]; Et += E[j]; printf("j=%d: terms=%llu  W_j=%.12Lf  E_j=%llu\n", j, cnt[j], W[j], E[j]); }
    ld D = 2 * Wt - 1 - 1.5L / lX;
    ld margin = A * D - 2.0L * (ld)Et;
    printf("X=%.6Le  log X=%.6Lf  A=X/logX=%.6Le  #good s=%zu\n", X, lX, A, good.size());
    printf("W=%.12Lf  2W=%.12Lf  needed 1+3/(2logX)=%.12Lf  D=%.12Lf\n", Wt, 2 * Wt, 1 + 1.5L / lX, D);
    printf("E=%llu  2E/A=%.12Lf\n", Et, 2.0L * (ld)Et / A);
    printf("margin A*D-2E = %.6Le  => %s\n", margin, (D > 0 && margin > 0) ? "TAIL CERTIFIED: pi(x) < 2T(x) for all x >= X" : "NOT CERTIFIED");
    return 0;
}
\end{lstlisting}

\subsection{Independent tail certificate (\texttt{tail\_check.py})}\label{code:tailpy}
\begin{lstlisting}[language=Python]
# Independent re-implementation of the tail certificate in Python.
# Template m = s * p_1 ... p_{j-1} (s 47-smooth, 2^j tau(s) | s, 53 <= p_1 < ... < p_{j-1} primes);
# included iff X > pi(p(m)) * m * log X, where p(m) = p_{j-1} (47 if j = 1).
import math, sys
import numpy as np
X = int(float(sys.argv[1])); PMAX = int(sys.argv[2])
logX = math.log(X); A = X / logX
sv = np.ones(PMAX + 1, dtype=bool); sv[:2] = False
for i in range(2, int(PMAX ** 0.5) + 1):
    if sv[i]: sv[i * i::i] = False
pr = np.nonzero(sv)[0]
q = [int(v) for v in pr[pr >= 53]]
NQ = len(q)
pis = [t + 16 for t in range(NQ)]                      # pi(q[t])
invpre = np.concatenate(([0.0], np.cumsum(1.0 / np.array(q, dtype=np.float64))))
pipre = [0]
for t in range(NQ): pipre.append(pipre[-1] + pis[t])
P = [2,3,5,7,11,13,17,19,23,29,31,37,41,43,47]
SMAX = int(A / 15) + 1
good = []
def dfs(i, s, tau):
    if i == 15:
        if s % tau == 0:
            r = s // tau; J = 0
            while r % 2 == 0: r //= 2; J += 1
            if J >= 1: good.append((s, J))
        return
    v, e = s, 0
    while True:
        dfs(i + 1, v, tau * (e + 1))
        if v * P[i] > SMAX: break
        v *= P[i]; e += 1
dfs(0, 1, 1)
def ok(m, pim): return X > pim * m * logX
def last_ok(base, lo):  # largest t >= lo with ok(base*q[t], pi(q[t])), or lo-1 if none (monotone in t)
    if not ok(base * q[lo], pis[lo]): return lo - 1
    a, b = lo, NQ - 1
    if ok(base * q[b], pis[b]): raise SystemExit("PMAX too small")
    while b - a > 1:
        mid = (a + b) // 2
        if ok(base * q[mid], pis[mid]): a = mid
        else: b = mid
    return a
W = [0.0] * 6; E = [0] * 6; nterm = [0] * 6
for (s, J) in good:
    if ok(s, 15): W[1] += 1.0 / s; E[1] += 15; nterm[1] += 1
    if J >= 2:
        c = last_ok(s, 0) + 1
        W[2] += invpre[c] / s; E[2] += pipre[c]; nterm[2] += c
    if J >= 3:
        for t1 in range(NQ - 1):
            c = last_ok(s * q[t1], t1 + 1) + 1
            if c == t1 + 1: break
            W[3] += (invpre[c] - invpre[t1 + 1]) / (s * q[t1]); E[3] += pipre[c] - pipre[t1 + 1]; nterm[3] += c - t1 - 1
    if J >= 4:
        for t1 in range(NQ - 2):
            any1 = False
            for t2 in range(t1 + 1, NQ - 1):
                c = last_ok(s * q[t1] * q[t2], t2 + 1) + 1
                if c == t2 + 1: break
                any1 = True
                W[4] += (invpre[c] - invpre[t2 + 1]) / (s * q[t1] * q[t2]); E[4] += pipre[c] - pipre[t2 + 1]; nterm[4] += c - t2 - 1
            if not any1: break
Wt = sum(W); Et = sum(E); D = 2 * Wt - 1 - 1.5 / logX
for j in range(1, 5): print(f"j={j}: terms={nterm[j]} W_j={W[j]:.12f} E_j={E[j]}")
print(f"#good s={len(good)} W={Wt:.12f} E={Et} D={D:.12f} 2E/A={2*Et/A:.12f} margin={A*D-2*Et:.6e}",
      "CERTIFIED" if D > 0 and A * D - 2 * Et > 0 else "NOT CERTIFIED")
\end{lstlisting}


\begin{thebibliography}{99}

\bibitem{Co99}
S.~Colton, \emph{Refactorable numbers --- a machine invention}, J. Integer Seq. \textbf{2} (1999), Article 99.1.2.

\bibitem{DR96}
M.~Deléglise and J.~Rivat, \emph{Computing $\pi(x)$: the Meissel, Lehmer, Lagarias, Miller, Odlyzko method}, Math. Comp. \textbf{65} (1996), 235--245.

\bibitem{FC}
Formal Conjectures (Google DeepMind and contributors), file \texttt{FormalConjectures/OEIS/111291.lean}, \url{https://github.com/google-deepmind/formal-conjectures}, commit \texttt{c094b97} (23 October 2026).

\bibitem{Go01}
X.~Gourdon, \emph{Computation of $\pi(x)$: improvements to the Meissel, Lehmer, Lagarias, Miller, Odlyzko, Deléglise and Rivat method}, preprint (2001).

\bibitem{KC90}
R.~E.~Kennedy and C.~N.~Cooper, \emph{Tau numbers, natural density, and Hardy and Wright's Theorem 437}, Internat. J. Math. Math. Sci. \textbf{13} (1990), 383--386.

\bibitem{OEIS}
OEIS Foundation Inc., \emph{The On-Line Encyclopedia of Integer Sequences}, sequences A033950 and A111291, \url{https://oeis.org}.

\bibitem{RS62}
J.~B.~Rosser and L.~Schoenfeld, \emph{Approximate formulas for some functions of prime numbers}, Illinois J. Math. \textbf{6} (1962), 64--94.

\bibitem{Sp85}
C.~Spiro, \emph{How often is the number of divisors of $n$ a divisor of $n$?}, J. Number Theory \textbf{21} (1985), 81--100.

\bibitem{primecount}
K.~Walisch, \emph{primecount} (version 7.10) and \emph{primesieve} (version 12.0), \url{https://github.com/kimwalisch/primecount}.

\bibitem{Ze02}
J.~Zelinsky, \emph{Tau numbers: a partial proof of a conjecture and other results}, J. Integer Seq. \textbf{5} (2002), Article 02.2.8.

\end{thebibliography}

\end{document}
